\section{Decision five: assembling the formulary}

Only now open the book. What follows is the whole variable surface of a Mass,
in the order it is used, with the rule for each slot and its locus. Everything
not listed here comes from the \work{Ordo Missae} unchanged.

\subsection{Where the text comes from}

Three books' worth of material meet in one Mass, and the code is explicit
about the seams.

\begin{itemize}
\item \textbf{The Ordinary and the Canon} supply everything invariable, and
the \work{Ordo Missae} itself carries the conditional directions --- the
alternative Gospel formula at the end, the two forms of the dismissal, the
alternative to the blessing.

\item \textbf{The Proper} supplies the formulary, either from the Proper of
the Season or from the Proper of Saints. A saint's entry in the Proper of
Saints is frequently not a complete Mass: it is a pointer to a Common plus
whatever is genuinely proper. St Patrick's entry on 17 March is typical ---
``Missa \latin{Statuit}, de Communi Confessoris Pontificis I loco, praeter
orationem sequentem,'' followed by his own collect, secret and postcommunion.
So the Mass is the Common's chants and readings with a proper set of orations.

\item \textbf{The Common} supplies the rest. When a feast is not in the Proper
of Saints at all, the whole Mass is taken from the Common (\rgm{305}\,b), and
where a Common offers several formularies the choice is the celebrant's. In
every Common, the epistles and gospels printed either within the individual
Masses or at the end of the whole Common ``sumi possunt in qualibet Missa de
eodem Communi'' --- they are interchangeable within that Common
(\rgm{305}\,b). For votive Masses of Our Lady, the Angels and the saints the
same route is prescribed by \rgm{315}, with a further instruction:
the parts that vary with the season and are missing from the formulary
``sumuntur e Communi festorum B. Mariae Virg.'' for Marian votives
(\rgm{309}), and in every case ``serventur \dots{} rubricae de mutandis
nonnullis partibus vel verbis, iuxta anni tempora et iuxta qualitatem mere
votivam huius Missae.''
\end{itemize}

The Missal performs those substitutions in print wherever the season demands
it. Its Masses carry directions such as ``Post Septuagesimam, omissis Alleluia
et versu sequenti, dicitur Tractus \dots{}'' and ``Tempore paschali antiphonae
ad Introitum et ad Offertorium sumuntur ex Missa \latin{Protexisti} \dots{}''
Where the book supplies no such direction for a season a Mass was never
expected to meet --- a votive of the Sacred Heart between Septuagesima and
Easter, for instance, whose feast Mass prints a Gradual and Alleluia and no
Tract --- the general rubrics do not fill the gap, and the slot is genuinely
open. Do not invent a Tract; consult the Ordo.

\subsection{The slot map}

\begin{slottable}
Prayers at the foot & Ordinary & Psalm \latin{Iudica me} with its antiphon and
the Confiteor with the absolution are said at every Mass, sung or read; they
are omitted with the following versicles and \latin{Aufer a nobis} and
\latin{Oramus te} in six named Masses (Purification after the blessing of
candles; Ash Wednesday after the imposition; Palm Sunday after the procession;
the Easter Vigil; the Rogation Mass after the Litany procession; certain Masses
following consecrations under the Pontifical). \rgm{424}\\
--- Psalm alone omitted & --- & The psalm alone is omitted in Masses
\latin{de Tempore} from Passion Sunday to Maundy Thursday, and in Masses of
the dead. \rgm{425}\\
Introit & Proper or Common & Antiphon, psalm verse, \latin{Gloria Patri},
antiphon repeated; absent only in the Easter Vigil Mass. \rgm{427}. The
\latin{Gloria Patri} is omitted in Masses \latin{de Tempore} from Passion
Sunday to Maundy Thursday and in Masses of the dead. \rgm{428}. In Paschaltide
a double Alleluia is added unless already there; outside Paschaltide Alleluia
is dropped from any Introit antiphon unless the Mass says otherwise.
\rgm{429}\\
Kyrie & Ordinary & Nine invocations, after the repeated Introit antiphon.
\rgm{430}\\
Gloria & Ordinary, conditional & Said when the Mass answers the Office of the
day and Te Deum was said at Matins; in the festive Masses of \rgm{302}; on
Maundy Thursday and at the Easter Vigil; in votive Masses of the first, second
and third class unless violet is used; and in fourth-class votives of the
Angels on any day and of Our Lady on Saturday. \rgm{431}. Omitted when Te Deum
is omitted at Matins, in every Mass using violet, in fourth-class votives
otherwise, and in Masses of the dead. \rgm{432}\\
Collect and other orations & Proper, Common, commemorations & See
\S\,6.3. \rgm{433}--\rgm{446}\\
Lessons before the Epistle & Proper & One lesson precedes the Epistle on Ember
Wednesdays, on Wednesday of the fourth week of Lent and on Wednesday of Holy
Week; five precede it on Ember Saturdays, with the shortening faculty of
\rgm{468} outside conventual and ordination Masses. \rgm{467}--\rgm{468}\\
Epistle & Proper or Common & \latin{Deo gratias} at the end. \rgm{466}\\
Gradual, Alleluia, Tract & Proper or Common & As printed in the Missal at its
own place. \rgm{469}\\
Sequence & Proper & Before the last Alleluia or after the Tract; \emph{omitted
in every votive Mass}; \latin{Dies irae} governed by \rgm{399}. \rgm{470}\\
Gospel & Proper or Common & \latin{Dominus vobiscum}, then \latin{Sequentia}
or \latin{Initium}. In Holy Week, before the Passion narrative, neither
\latin{Dominus vobiscum} nor \latin{Sequentia} nor \latin{Gloria tibi} is
said, and \latin{Laus tibi, Christe} is not answered at the end.
\rgm{471}--\rgm{472}\\
Homily & --- & After the Gospel, especially on Sundays and days of precept, a
short homily is to be given as opportunity allows; if given by a priest other
than the celebrant, the Mass is suspended and resumed only when it is
finished. \rgm{474}\\
Creed & Ordinary, conditional & See \S\,6.4. \rgm{475}--\rgm{476}\\
Offertory antiphon & Proper or Common & After the Creed, or after the Gospel
or homily if there is none; absent only in the Easter Vigil Mass; Alleluia
added in Paschaltide unless already there, and a final Alleluia already
present is kept outside Paschaltide except from Septuagesima to Easter.
\rgm{477}--\rgm{478}\\
Secrets & Same sources as the collects & Said silently, without
\latin{Dominus vobiscum} and \latin{Oremus}; as many as there were orations at
the beginning of Mass, in the same order and with the same conclusions; the
conclusion of the last is silent until \latin{Per omnia saecula saeculorum}.
\rgm{480}--\rgm{481}\\
Preface & Proper, season, or common & See Section~7. \rgm{482}--\rgm{499}\\
Canon & Ordinary & Said silently. Proper \latin{Communicantes}, \latin{Hanc
igitur} and \latin{Qui pridie} are noted in the Masses that have them; within
the octaves of Christmas, Easter and Pentecost the proper
\latin{Communicantes} and \latin{Hanc igitur} are said even in Masses that are
not of the octave, ``etiamsi praefatione propria gaudeant.''
\rgm{500}--\rgm{501}\\
Communion antiphon & Proper or Common & Alleluia added in Paschaltide unless
already there; a final Alleluia already present is kept outside Paschaltide
except from Septuagesima to Easter. \rgm{504}\\
Postcommunions & Same sources as the collects & ``eodem numero, modo et ordine
ac orationes in principio Missae.'' \rgm{505}\\
Prayer over the people & Proper & Added after the last postcommunion in Masses
\emph{of the ferias} of Lent and Passiontide, the Sacred Triduum excepted;
always under its own conclusion, preceded by \latin{Oremus. Humiliate capita
vestra Deo}; said even when three postcommunions have already preceded.
\rgm{506}\\
Dismissal & Ordinary, conditional & \latin{Ite, missa est}; but
\latin{Benedicamus Domino} at the evening Mass of Maundy Thursday followed by
the solemn reposition and in other Masses followed by a procession; a double
Alleluia added within the octave of Easter in Masses \latin{de Tempore};
\latin{Requiescant in pace} in Masses of the dead, answered \latin{Amen}.
\rgm{507}\\
Blessing & Ordinary, conditional & Given after \latin{Placeat}; omitted only
where \latin{Benedicamus Domino} or \latin{Requiescant in pace} was said.
\rgm{508}\\
Last Gospel & Ordinary or Proper & See Section~7. \rgm{509}--\rgm{510}\\
\end{slottable}

\subsection{The orations: number, order and conclusion}

This is the most intricate part of the assembly and the part in which a hand
missal from before 1955 will actively mislead.

\rgm{433} defines what an oration is for counting purposes: the oration of the
Mass being celebrated; the orations of a commemorated Office and of any
occurring commemoration; other orations prescribed by the rubrics; the prayer
imposed by the local Ordinary; and the votive prayer that may be added at the
celebrant's choice on certain days. \rgm{434} then makes the ceiling apply to
all of them together:

\begin{rulebox}{\rgm{434}}
Numero orationum pro singulis diebus liturgicis statuto complectuntur tam
oratio Missae et commemorationes quam aliae orationes sive a rubricis
praescriptae sive ab Ordinario imperatae sive votivae. Proinde, post orationem
Missae: a) in diebus liturgicis I classis, in Missis votivis I classis, et in
Missis in cantu non conventualibus, nulla alia admittitur oratio, praeter
orationem sub unica conclusione dicendam et unam commemorationem
privilegiatam, salvo praescripto n.~333; b) in dominicis II classis, nulla
alia admittitur oratio, praeter commemorationem festi II classis, quae tamen
omittitur si commemoratio privilegiata facienda sit; c) in aliis diebus
liturgicis II classis et in Missis votivis II classis una tantum alia
admittitur oratio, scilicet aut una privilegiata aut una ordinaria; d) in
diebus liturgicis III et IV classis et in Missis votivis III et IV classis duae
tantum admittuntur orationes.
\end{rulebox}

\begin{rulebox}{\rgm{435}}
Quaelibet oratio, quae numerum pro singulis diebus liturgicis statutum superet,
omittitur; profecto numerum ternarium orationum nullo praetextu excedere
licet.
\end{rulebox}

Three practical consequences:

\begin{substeps}
\item \textbf{The sung non-conventual Mass is treated as a first-class day for
this purpose.} \rgm{434}\,a puts ``Missis in cantu non conventualibus'' in the
same clause as first-class days: one privileged commemoration, and no more,
whatever the calendar says. A parish sung Mass on a fourth-class Tuesday
therefore has fewer orations than the Low Mass said at the side altar an hour
earlier. \rg{108} is the reason: ordinary commemorations are made only at
Lauds, at the conventual Mass, and at Low Masses.

\item \textbf{The absolute cap is three, always.} \latin{Nullo praetextu}.

\item \textbf{Under one conclusion, or under a second.} The Mass's own oration
is always said under its own conclusion unless another is to be joined to it
(\rgm{436}); commemorations, the imposed prayer and the votive prayer are
always said under a second conclusion (\rgm{437}). Only three kinds of oration
may be joined under a single conclusion --- a ritual prayer, the prayer of an
impeded first- or second-class votive Mass, and a prayer the rubrics expressly
direct to be so said (\rgm{444}) --- and only one may be joined, the rest being
dropped in the order \rgm{444} lists (\rgm{445}). An oration under a single
conclusion counts as one with the Mass's own, and ``dicenda est etiam in Missis
in cantu'' (\rgm{446}): it survives the sung-Mass restriction.
\end{substeps}

Two smaller rules save trouble. If two orations are composed in nearly the same
words in their first or second part, the later one is changed --- to the
following Sunday's or feria's if it is of the season, to another of the same
or a similar Common if it is of a saint --- and simply omitted if it is the
imposed prayer (\rgm{438}). And in the orations of a transferred or
repositioned Office the words \latin{hanc}, \latin{hodiernam},
\latin{praesentem diem} are not to be altered (\rgm{439}) --- ``this day'' is
said on the day the feast has been given.

The imposed prayer deserves its own note because it is so often assumed to be
available. \rgm{454} allows the local Ordinary to impose a prayer ``occurrente
gravi et publica necessitate aut calamitate.'' It may be one only; it must be
said by all priests celebrating in the diocese, even the exempt; it is never
said under a single conclusion with the Mass's oration but after the privileged
commemorations; and

\begin{rulebox}{\rgm{457}\,d}
prohibetur omnibus diebus liturgicis I et II classis, in Missis votivis I et II
classis, in Missis in cantu et quoties commemorationes privilegiatae numerum
pro singulis diebus liturgicis statutum compleverint.
\end{rulebox}

\noindent
\latin{In Missis in cantu} --- in every sung Mass, without the qualification
``non conventualibus'' that \rgm{434}\,a carries. The votive prayer of
\rgm{461} is narrower still: any priest may add one at his discretion ``in
omnibus Missis lectis non conventualibus diebus liturgicis IV classis'' --- read
Masses, not conventual, fourth-class days only --- and it is placed last
(\rgm{463}).

\subsection{The Creed}

\begin{rulebox}{\rgm{475}}
Post Evangelium aut homiliam, dicitur symbolum: a) in qualibet dominica, etsi
eius Officium alicui festo locum cedat, vel Missa votiva II classis
celebretur; b) in festis I classis et in Missis votivis I classis; c) in
festis II classis Domini et B. Mariae Virg.; d) per octavas Nativitatis
Domini, Paschatis et Pentecostes, etiam in festis occurrentibus et in Missis
votivis; e) in festis nataliciis Apostolorum et Evangelistarum, necnon in
festis Cathedrae S. Petri et S. Barnabae Ap.
\end{rulebox}

Clause (a) is the one that surprises. The Creed belongs to the \emph{day},
not to the formulary: on any Sunday it is said even when the Sunday's Office
has yielded to a feast, and even when a second-class votive Mass is being said.
\rgm{476} then excludes it in the chrism and Maundy Thursday Masses and at the
Easter Vigil; in second-class feasts other than those just named; in
second-class votive Masses; in festive and votive Masses of the third and
fourth class; ``ratione alicuius commemorationis in Missa occurrentis'' --- a
commemoration never brings the Creed with it --- and in Masses of the dead.

\subsection{Colour}

The five colours are white, red, green, violet and black (\rg{117}), with
legitimate indults and customs preserved and a faculty for episcopal
conferences in mission regions to substitute a more suitable colour where a
Roman colour's meaning conflicts with an indigenous tradition, ``hoc tamen non
fiat inconsulta S. Rituum Congregatione.'' Colour follows the Office and Mass
of the day or the other Mass being celebrated (\rg{117}), and the detailed
assignments occupy \rg{119}--\rg{132}: rose is permitted on the Third Sunday
of Advent and the Fourth of Lent, ``sed in Officio et Missa diei dominici
tantum'' (\rg{131}).

One exception is worth carrying: in fourth-class read votive Masses that are
not conventual, the colour of the Office of the day may be used instead of the
Mass's own, ``servato tamen colore violaceo et nigro unice pro Missis quibus
per se competit'' (\rgm{323}, referred forward by \rg{118}). And because
\rgm{431}\,d and \rgm{432}\,b make the Gloria depend on whether violet is
worn, the colour decision can silently change the Ordinary.
